\documentclass[12pt]{article}
\usepackage[utf8]{inputenc}
\usepackage[T1]{fontenc}
\usepackage{hyperref}
\title{Parsi-Dari Language Resources for K-12 Curriculum Development}
\author{Compiled by Claude Sonnet 5.5 (claude-sonnet-5-5) for Andr\'as Kornai}
\date{}
\begin{document}
\maketitle

\section{Introduction}
The code prd (Parsi-Dari) is a disputed entry. Glottolog lists it only as a bookkeeping entry and treats it as a duplicate of Zoroastrian Dari (ISO 639-3 gbz), the Dari variety of the Zoroastrian community around Yazd and Kerman in Iran. Searches found no resources filed specifically under ``Parsi-Dari''. The resources below concern Zoroastrian Dari (also called Behdin\={a}n\={\i} or Dari; ``Gabri'' is an outsider term that speakers resent). They are listed because they are the closest material found, and they should not be assumed to describe a separate language. The Kermani variety is reported to be extinct, and the Yazdi variety is still spoken.

\section{Primary Educational Resources}
No school curriculum, primer or classroom material was found in search results. The closest item is a learner's book on the Kerman variety:
\begin{itemize}
\item Gholami and Farahmand, \emph{Zoroastrian Dari (Behdini) in Kerman} (2016), a learner's textbook with grammar, vocabulary and texts. Information page: \url{https://www.biblioiranica.info/zoroastrian-dari-behdini-in-kerman/}
\end{itemize}

\section{Linguistic and Reference Materials}
\begin{itemize}
\item Glottolog bookkeeping entry for Parsi-Dari: \url{https://glottolog.org/resource/languoid/iso/prd}
\item Wikipedia, Zoroastrian Dari language: \url{https://en.wikipedia.org/wiki/Zoroastrian_Dari_language}
\item Encyclop\ae dia Iranica, Behdinan dialect: \url{https://www.iranicaonline.org/articles/behdinan-dialect}
\item Cambridge, Iranian Studies article on remnants of Zoroastrian Dari in manuscript colophons: \url{https://resolve.cambridge.org/core/journals/iranian-studies/article/remnants-of-zoroastrian-dari-in-the-colophons-and-salmargs-of-iranian-avestan-manuscripts/1A56E34168FCE371B4608E53F08B686E}
\end{itemize}
A Persian-language dictionary, \emph{V\={a}zhen\={a}me-ye guyesh-e behdin\={a}n-e shahr-e Yazd}, vol.\ I (Tehran, 1995), appears in a bibliography. No online copy was found.

\section{Religious and Community Resources}
No community organization or scripture translation in the language was identified in the search results.

\section{Digital and Audio Resources}
No audio, video or software resources were found.

\section{Curriculum Development Context}
Because the code is disputed and almost nothing is published under it, any curriculum would in practice rest on Zoroastrian Dari documentation, and would need work with Yazd and Kerman community members. The language is primarily oral, so written material is scarce.

\end{document}